\chapter{Classification of Integral Affine Closed Surfaces}
As we have seen in the previous chapter, the integral affine stucture induced
by a regular Lagrangian torus fibration is an invariant of the base space of
the fibration. Moreover, given an integral affine structure on a manifold $B$,
we can construct a Lagrangian fibration by taking the quotient of $T^*B$ by the
lattice inducing the integral affine structure. We shall see in Chapter
\ref{ch:classify-fibr} that after having fixed an integral affine structure on
$B$, all Lagrangian fibrations over $B$ will locally look like $T^*B/\Lambda
    \to B$ where $\Lambda \subset T^*B$ is the period lattice inducing the integral
affine structure.

In order to classify the regular Lagranigan torus fibrations over a manifold
$B$, the first step is to classify the integral affine structures on $B$. The
goal of this chapter is to carry this out for compact surfaces.

\section{Classification of integral affine circles}
It is easy to classify the integral affine structures on the circle $S^1$ and
so for completeness we will also give the classification of these.

The fundamental group of $S^1$ is the integers $\Z$ and hence by Theorem
\ref{thm:complete-parallel-volume} any affine structure will be complete if it
possesses a parallel volume form. The holonomy representation $\hol :
    \pi_1(S^1) \to \affz(\R)$ is determined by the image of $1 \in \Z$ and so the
possible affine holonomy groups are $$ A = \langle (1, a) \rangle \quad
    \text{or} \quad B =\langle(-1, b) \rangle$$ with $a,b \in \R$. We note that $B
    \cong \Z/2\Z$ and in fact, this affine structure would be radiant as
$\frac{b}{2}$ is a fixed point. However, Corollary \ref{corr:radiant-expansion}
says that the linear holonomy would contain an expansion but all the
eigenvalues of the elements of $B$ have modulus one. Hence $B$ cannot occur as
the affine holonomy group of an integral affine structure over $S^1$.

Now, since the affine holonomy group $A$ preserves a parallel volume form, by
Theorem \ref{thm:complete-parallel-volume}, this affine structure is complete.
The integral affine structures on $S^1$ are therefore of the form $S^1 = \R /
    A$. By Lemma \ref{lemma:equivalent-dev-pair}, complete integral affine
structures are equivalent when their affine holonomy groups are conjugate. The
following lemma classifies the integral affine structures on $S^1$.

\begin{lemma}\label{lemma:classification-integral-affine-circles}
    The integral affine structures on $S^1$ are determined by a positive real number $a>0$ and are of the form $\R/A$ where
    $$
        A = \langle (1, a) \rangle.
    $$
\end{lemma}
\begin{proof}
    Our previous discussion showed that all integral affine structures on $S^1$ are of the form $\R/A$ where the translational component is any real number. It remains to check when two structures are equivalent.
    We note that the only possible linear affine holonomy is $\{\pm 1\}$ and that if it is $1$ conjugating $A$ leaves it unchanged. We therefore compute the conjugation by an element with negative linear holonomy.
    Letting $b \in \R$ we compute
    \begin{align*}
        (-1, b)(1,a)(-1,b)^{-1} & = (-1, b)(1,a)(-1,b) \\
                                & = (1,-a)
    \end{align*}
    but this generates the same affine holonomy group as $(1,a)$ completing the proof since if $a \neq 0$ the action is clearly free.
\end{proof}

\section{Some basic results}
We now move on to classifying the integral affine structures over compact
surfaces which requires a bit more work. We first discuss our approach to the
classification problem and state some important results in the theory of affine
manifolds. The first task is to determine the possible compact surfaces that
can admit an affine structure. The following theorem classifies the compact,
connected $2$-manifolds.

\begin{theorem}
    Every compact, connected $2$-manifold is homeomorphic to the sphere, the connected sum of tori or the connected sum of $\mathbb{R}P^2$.
\end{theorem}
\begin{proof}
    This can be found in \cite{Hatcher}.
\end{proof}

We make use of the following theorem from Milnor relating to the existence of a
flat connection on closed orientable surfaces.
\begin{theorem}[Milnor \cite{Milnor1957/58}]
    Let $M^2$ be a closed orientable surface of genus $g$. Then if $g \geq 2$ it does not possess a connection with zero curvature.
\end{theorem}

\begin{corollary}
    The possible orientable closed surfaces admitting an integral affine structure are the sphere $S^2$ or the torus $\mathbb{T}^2$.
\end{corollary}

\begin{proof}
    This follows immediately from the classification of closed 2-surfaces and the previous theorem.
\end{proof}

However, the following simple argument shows that the sphere $S^2$ does not
admit such a structure.

\begin{prop}
    The sphere $S^2$ does not admit an affine structure.
\end{prop}

\begin{proof}
    Suppose there was such an affine structure. Then
    $$\dev : S^2 \to \mathbb{R}^2$$
    is a local diffeomorphism and hence an open map. However $\dev$ is also continuous, and so $\dev(S^2)$ is a compact subset of $\mathbb{R}^2$ and hence closed.
    Now, as $\dev(S^2) \neq \emptyset$, by connectivity $\dev(S^2) = \mathbb{R}^2$ but $\mathbb{R}^2$ is not compact.
\end{proof}
We have the following immediate corollary.
\begin{corollary}
    The only closed, orientable surface admitting an (integral) affine structure is $\mathbb{T}^2$.
\end{corollary}

We will deal with the non-orientable case by passing to the orientable double
cover.
\begin{lemma}
    Every closed, non-orientable surface has a double cover that is orientable. Furthermore, the genus of the double cover is given by
    $$g = \bar{g} -1 $$
    where $\bar{g}$ is the number of connected sums of $\mathbb{R}P^2$.
\end{lemma}
\begin{proof}
    This is proved in \cite{Hatcher}.
\end{proof}

This lemma leads to the following corollary.
\begin{corollary}
    The only closed surfaces admitting an (integral) affine structure are the torus $\mathbb{T}^2$ and the Klein bottle $\mathbb{K}^2$.
\end{corollary}
\begin{proof}
    As any covering space of an integral affine structure inherits an integral affine structure by (TODO: do this in (G,X)-structure section), it follows that the
    orientable double cover of a non-orientable integral affine surface must be $\mathbb{T}^2$. Furthermore, it must have non-orientable genus $\bar{g}= 1 \text{ or } 2$. If $\bar{g}=1$ our surface is homeomorphic to $\mathbb{R}P^2$ which is covered by
    $S^2$ and hence admits no such structure. If $\bar{g} = 2$ we get the Klein bottle which is covered by $\mathbb{T}^2$ and hence it can admit such a structure.
\end{proof}

\section{Classification of Integral Affine structures on the torus}
\subsection{Completeness and freeness of the action}
We now look at how we can classify the integral affine structures on the torus.
The following observation will greatly simplify the classification process.

\begin{lemma}
    Any integral affine structure on the torus is complete.
\end{lemma}
\begin{proof}
    We have that the fundamental group $\pi_1(\mathbb{T}^2) = \mathbb{Z} \times \mathbb{Z}$ is abelian and hence nilpotent. As the torus is orientable,
    our linear holonomy group $\Gamma \subseteq \text{SL}(2, \mathbb{Z})$ by Corollary \ref{lemma:volume} and now by Theorem \ref{thm:complete-parallel-volume} this is complete.
\end{proof}

(TODO: This might be a good place for discussion on completeness stuff. Auslander-Markus show that every complete affine manifold is quotient by linear holonomy).

From the above discussion and (TODO: recap relevant results), classifying
complete integral affine structures on a surface $M$ amounts to classifying
injective homomorphisms from the fundamental group into the group
$\text{Aff}_{\mathbb{Z}}(\mathbb{R}^2)$ which induce a free and proper action
on $\mathbb{R}^2$.

\begin{lemma}
    If $g=(A,b) \in \affz(\mathbb{R}^2)$ acts freely on $\mathbb{R}^2$, then
    $$\text{if } \det(A) =1 \text{ then } \text{Tr}(A) = 2$$ and,
    $$\text{if } \det(A) =-1 \text{ then } \text{Tr}(A) = -1.$$\label{lemma:traces-free-action}
\end{lemma}
\begin{proof}
    Since the action by $g$ is free, this implies
    $$Ax +b =x$$
    has no solution. But now,
    \begin{align*}
         & (A-I)x = -b
    \end{align*}
    $\text{has no solutions for all } x\in \mathbb{R}^2$ and so $\det(A-I) = 0$. Calculating this determinant yields
    $$1- \text{Tr}(A) + \det(A) = 0$$
    and substituting in the appropriate values of $\det(A)$ gives the required results.

\end{proof}

We note also that $1$ is an eigenvalue of such an $A$ as by the above proof
$\det(A - 1\cdot I) = 0$. Using this result we can find a nice basis to express
the linear part of our element $g \in \affz(\mathbb{R}^2)$.

\begin{lemma}
    Let $g= (A,b) \in \affz(\mathbb{R}^2)$ act freely. Then there is an integral basis of $\mathbb{R}^2$ such that
    $$A = \begin{pmatrix}
            1 & n     \\
            0 & \pm 1
        \end{pmatrix}$$ \label{lemma:torus-matrixform}
\end{lemma}
\begin{proof}
    By our previous lemma, $A$ must have $1$ as an eigenvalue. We can take an eigenvector $e_1$ such that $e_1$ is integral, i.e. its components are integers.
    Now as $\text{Tr}(A) = 2 \text{ or } 0$, it follows that we can pick our second basis vector such that
    $$A = \begin{pmatrix}
            1 & n     \\
            0 & \pm 1
        \end{pmatrix}$$
    since we need $A \in \text{GL}(2, \mathbb{Z})$. Furthermore, by changing orientation if necessary we may assume $n \geq 0$.
\end{proof}

\begin{lemma}
    \label{lemma:skew-matrix-form}

    Let $H \subset \affz(\mathbb{R}^2)^+$ be a subset of the elements of
    determinant one and suppose $H$ acts freely on $\mathbb{R}^2$. Then there is an
    integral basis such that for all $(A,b)\in H$ we can write $$A = \begin{pmatrix}
            1 & n \\
            0 & 1
        \end{pmatrix}$$
    for $n \in \mathbb{Z}.$
\end{lemma}

\begin{proof}
    Let $(A, a),(B,b) \in H$. Now by Lemma \ref{lemma:torus-matrixform}, we can write
    $$A = \begin{pmatrix}
            1 & n \\
            0 & 1
        \end{pmatrix}, \quad B = \begin{pmatrix}
            a & b \\
            c & d
        \end{pmatrix}$$
    with $n \neq 0$. Now by Lemma \ref{lemma:traces-free-action}, we have that
    \begin{align*}
        \text{Tr}(A^kB) & = a + nck + d        \\
                        & = \text{Tr}(B) + nck \\
                        & = 2.
    \end{align*}
    As $n\neq 0$ and this holds for all $k \in \mathbb{Z}$, it follows that $c=0$.

    But now, $\det B = ad = 1$ and $\text{Tr}(B) = a+d = 2$ and so $a=d=1$ giving
    the required result as $b \in \mathbb{Z}$.
\end{proof}

We have the following corollary of this result.
\begin{corollary}
    \label{corr:rel-prime}
    For any two elements $(A,a), (B,b) \in H$, there exists relatively prime $(k,l) \neq (0,0)$ such that $A^kB^l = 1$.
\end{corollary}
\begin{proof}
    We simply take $k = \frac{m}{(m,n)}$ and $l = \frac{-n}{(m,n)}$ and then
    $$\begin{pmatrix}
            1 & n \\
            0 & 1
        \end{pmatrix}^k\begin{pmatrix}
            1 & m \\
            0 & 1
        \end{pmatrix}^l = 1$$.
\end{proof}

\subsection{Finding the possible generators}
We have now built up the required machinery to begin the classification of the
integral affine structures on the torus $\mathbb{T}^2$.

\begin{theorem}\label{thm:classification-torus}
    Any complete integral affine structure on $\mathbb{T}^2$ has its affine holonomy group generated by one of the following subgroups of $\affz(\mathbb{R}^2)$
    \begin{enumerate}
        \item $\mathcal{A}_{a,b,c,d}$ - These integral affine structures are the quotient of $\mathbb{R}^2$ by a lattice of pure translations. The generators in $\affz(\mathbb{R}^2)$ are given by
              $$\left( 1, \begin{pmatrix}
                          a \\ c
                      \end{pmatrix} \right) \text{ and } \left( 1, \begin{pmatrix}
                          b \\ d
                      \end{pmatrix} \right)$$
              with $a,b,c,d \in \mathbb{R}$ and $ad-bc \neq 0$.
        \item $\mathcal{B}_{n,a,b}$ - The generators in $\affz(\mathbb{R}^2)$ are given by
              $$\left( 1, \begin{pmatrix}
                          a \\ 0
                      \end{pmatrix} \right) \text{ and } \left( \begin{pmatrix}
                          1 & n \\
                          0 & 1
                      \end{pmatrix}, \begin{pmatrix}
                          0 \\ b
                      \end{pmatrix} \right)$$
              with $a,b > 0$ and $n \in \mathbb{N}$.
    \end{enumerate}
\end{theorem}
\begin{proof}
    Let $\pi_1(\mathbb{T}^2) = \langle \tilde{C}, \tilde{D} \rangle$ and denote by $C,D$ the images of these generators under the affine holonomy representation.
    We wish to show that these generators reduce to the form of either $\mathcal{A}_{a,b,c,d}$ or $\mathcal{B}_{n,a,b}$. We can do this by
    \begin{enumerate}
        \item Changing the generators $\tilde{C}, \tilde{D}$ of the $\pi_1(\mathbb{T}^2)$.
        \item Conjugating the images of these generators by some $g \in \affz^+(\mathbb{R})$
              - this corresponds to pre-composing the developing map by $g$ which preserves
              the affine structure.
    \end{enumerate}
    If both $C= \left( 1 ,\begin{pmatrix}
                a \\ c
            \end{pmatrix}\right)$ and $D =\left(1, \begin{pmatrix}
                b \\ d
            \end{pmatrix}\right)$ are pure translations, e.g. the linear component is the identity, injectivity forces that
    our translational vectors are linearly independent, otherwise $\hol(k\tilde{C} + m\tilde{D}) = 0$ for some $k,m\neq 0 \in \mathbb{R}$. In this
    case our lattice in $\mathbb{R}^2$ would be one-dimensional. This gives the condition that $ad-bc \neq 0$ and completes the first class of lattices. \\

    Now suppose that $C,D$ are not pure translations. Then by Corollary
    \ref{corr:rel-prime}, there exists relatively prime $k,l \in \mathbb{Z}$ such
    that $C^kD^l$ is a translation. We can therefore change generators in
    $\mathbb{Z}^2$ to $\tilde{A}=\tilde{C}^k\tilde{D}^l$ and $\tilde{B}= \begin{pmatrix}
            a \\ b
        \end{pmatrix}$. In the basis $\tilde{C}, \tilde{D}$ we have that the matrix
    $$M = \begin{pmatrix}
            k & a \\
            l & b
        \end{pmatrix}$$
    has $\det M = kb - la$ and since $k,l$ are relatively prime by Bezout's theorem we can always pick $a,b \in \mathbb{Z}$ such that
    this determinant is 1 and hence an automorphism of our lattice.

    We denote the images of our new generators under the affine holonomy
    representation by $$A' = \left( I, a =\begin{pmatrix}
                a_1 \\ a_2
            \end{pmatrix} \right) \text{ and } B' = \left( B, b =\begin{pmatrix}
                b_1 \\ b_2
            \end{pmatrix} \right).$$
    Now by Lemma \ref{lemma:skew-matrix-form}, we can find a basis of $\R^2$ such
    that $B = \begin{pmatrix}
            1 & n \\
            0 & 1
        \end{pmatrix}$ with $n>0$. Now $\hol(\pi_1(\mathbb{T}^2))$ is abelian so we require that $A'B'= B'A'$ which gives the condition
    $$
        Ba = a
    $$
    this implies that
    \begin{align*}
        \begin{pmatrix}
            1 & n \\
            0 & 1
        \end{pmatrix}\begin{pmatrix}
                         a_1 \\ a_2
                     \end{pmatrix} & = \begin{pmatrix}
                                           a_1 + na_2 \\
                                           a_2
                                       \end{pmatrix}   \\
                                     & = \begin{pmatrix}
                                             a_1 \\
                                             a_2
                                         \end{pmatrix}
    \end{align*}
    which gives $a_2=0$ as $n>0$ and so $a = \begin{pmatrix}
            a_1 \\
            0
        \end{pmatrix}$. We note that $a_1 \neq 0$ as otherwise the action is not free. For freeness, we also require that $A^lB^k$ acts freely for all $k,l \in \mathbb{Z}$. We check
    \begin{align*}
        \begin{pmatrix}
            1 & n \\
            0 & 1
        \end{pmatrix}^k \begin{pmatrix}
                            x \\ y
                        \end{pmatrix} + \begin{pmatrix}
                                            kb_1 + la_1 \\ kb_2
                                        \end{pmatrix} & =
        \begin{pmatrix}
            x + kny +kb_1 + la_1 \\ y + kb_2
        \end{pmatrix}                    \\
                                        & = \begin{pmatrix}
                                                x \\ y
                                            \end{pmatrix}
    \end{align*}
    If $b_2 = 0$ then there is a fixed point for $y = \dfrac{-kb_1-la_1}{kn} \in \mathbb{R}$ so we must have $b_2 \neq 0$. Conversely, if $b_2 \neq 0$ and the action is not free,
    this requires $k=0$ and so $x+la_1 = a_1$ for $l \neq 0$ but $a_1 \neq 0$ and hence there are no solutions. Now, to show that this takes the form of $\mathcal{B}_{n,a,b}$, we conjugate by
    $$T = \left( I, \begin{pmatrix}
                0 \\ \frac{-b_1}{n}
            \end{pmatrix}\right)$$
    this fixes $A'$ and so we compute $$T^{-1}B'T = \left(B, \begin{pmatrix}
                0 \\ b_2
            \end{pmatrix} \right)$$
    and so we may assume $b_1 =0$. We can also make the change $A' \mapsto A'^{-1}$ and $B' \mapsto B'^{-1}$ and assume $a_1,b_2 >0$ which completes the proof.
\end{proof}

\subsection{When are two structures equivalent?}

This gives the possible integral affine structures on the 2-torus. However we
still have to tackle the question of when $2$ lattices give the same integral
affine structure. This occurs when the subgroups of the integral affine group
determined by the holonomy representation are conjugate to eachother.

We have the following classification result.
\begin{lemma}
    Lattices from different series do not induce equivalent integral affine structures.
\end{lemma}
\begin{proof}
    Lattices in series $\mathcal{A}_{a,b,c,d}$ are generated by pure translations and hence cannot be conjugated into lattices of the form $\mathcal{B}_{n,a,b}$.

\end{proof}

\begin{theorem}\label{thm:torus-series-a}
    Lattices $\mathcal{A}_{a,b,c,d}$ and $\mathcal{A}_{a',b',c',d'}$ induce
    equivalent integral affine structures if and only if $$X = GX'H$$ where $X = \begin{pmatrix}
            a & b \\
            c & d
        \end{pmatrix}$ and $X' = \begin{pmatrix}
            a' & b' \\
            c' & d'
        \end{pmatrix}$ with $G,H \in GL(2,\mathbb{Z}).$
\end{theorem}
\begin{proof}
    Performing lattice automorphisms in $\pi_1(\mathbb{T}^2)$ and conjugating in
    the image of the generators in $\affz(\R^2)$ preserves the integral affine
    structure.

    Let $\pi_1(\mathbb{T}^2) = \langle v_1, v_2\rangle$ with $\hol(v_1) = \left(I, \begin{pmatrix}
                a \\ c
            \end{pmatrix} \right)$ and $\hol(v_2) = \left(I, \begin{pmatrix}
                b \\ d
            \end{pmatrix} \right)$ and set $X= \begin{pmatrix}
            a & b \\
            c & d
        \end{pmatrix}$. Let $H \in GL(2,
        \mathbb{Z})$ and $g = (G, w) \in \affz(\mathbb{R}^2)$. We note that conjugating a pure translation $\left(I, v\right)$ by $g$ gives the element $\left(I, Gv\right)$. We compute
    $$
        \hol(Hv_1) = \left(I, XH\begin{pmatrix}
                1 \\ 0
            \end{pmatrix}\right)
    $$
    and
    $$
        \hol(Hv_2) = \left(I, XH\begin{pmatrix}
                0 \\ 1
            \end{pmatrix}\right).
    $$
    This corresponds to an automorphism of our lattice in $\pi_1(\mathbb{T}^2)$. We can also conjugate generators by $g = (G, w) \in \affz(\mathbb{R}^2)$ which gives
    $$
        g\circ \hol(Hv_1) \circ g^{-1} = \left(I, GXH\begin{pmatrix}
                1 \\ 0
            \end{pmatrix}\right)
    $$
    and
    $$
        g\circ \hol(Hv_2) \circ g^{-1} = \left(I, GXH\begin{pmatrix}
                0 \\ 1
            \end{pmatrix}\right)
    $$
    and so it follows that the integral affine structures are equivalent if and only if $X' = GXH$. TODO: finish this proof.
\end{proof}

\begin{theorem}
    Lattices in $\mathcal{B}_{n,a,b}$ never induce the same integral affine
    structure.
\end{theorem}
\begin{proof}
    Suppose that $\mathcal{B}_{n,a,b} =\langle\left( 1, \bar{a} \right),  \left( B, \bar{b} \right) \rangle$ and $\mathcal{B}_{n',a',b'}= \langle\left( 1, \bar{a}' \right),  \left( B', \bar{b}' \right) \rangle$ induce equivalent integral affine structures. We must have that
    \begin{equation}\label{eq:first-conjugacy}
        g\tilde{A}g^{-1} = \tilde{A}'^k\tilde{B}'^r
    \end{equation}
    \begin{equation}
        \label{eq:second-conjugacy}
        g\tilde{B}g^{-1} = \tilde{A}'^l\tilde{B}'^s
    \end{equation}
    where $\tilde{A}, \tilde{B}$ and $\tilde{A}', \tilde{B}'$ are the images of the generators of $\pi_1(\mathbb{T}^2)$ of the structures $\mathcal{B}_{n,a,b}$ and $\mathcal{B}_{n',a',b'}$ and $g = (G = [g_{ij}], \bar{h})\in \affz(\R^2)$. The matrix
    $$M = \begin{pmatrix}
            k & l \\ r & s
        \end{pmatrix} \in GL(2, \mathbb{Z})$$
    corresponds to an automorphism of our lattice. This implies that
    \begin{align*}
        GBG^{-1} = \begin{pmatrix}
                       1 & * \\
                       0 & 1
                   \end{pmatrix}
    \end{align*}
    where $B = \begin{pmatrix}
            1 & n \\
            0 & 1
        \end{pmatrix}$. Computing this condition yields
    $$
        \begin{pmatrix}
            *                                       & * \\
            g_{21}g_{22} - ng_{21}^2 - g_{21}g_{22} & *
        \end{pmatrix} = \begin{pmatrix}
            1 & * \\
            0 & 1
        \end{pmatrix}
    $$
    and so $g_{21}g_{22} - ng_{21}^2 - g_{21}g_{22} = 0$ which gives that $g_{21} = 0$.

    Now, inspecting the linear part of $g\tilde{A}g^{-1}$ as the linear part of
    $\tilde{A}$ is the identity this must also be a translation i.e. $\tilde{A}'^k
        \tilde{B}'^r$ is a translation. Again, inspecting the linear components we must
    have that $r=0$. Our matrix $M$ corresponding to the automorphism of our
    lattice now looks like $$M = \begin{pmatrix}
            k & l \\
            0 & s
        \end{pmatrix}$$
    and so $k,s = \pm 1$ as this matrix needs to have determinant $\pm 1$. Now from Equation $\ref{eq:first-conjugacy}$ we obtain
    \begin{align*}
        \left(G, \bar{h}\right) \left(I, \begin{pmatrix}
                                                 a \\ 0
                                             \end{pmatrix}\right) \left(G^{-1}, -G^{-1}\bar{h}\right) & = \left(I, G\begin{pmatrix}
                                                                                                                        a \\ 0
                                                                                                                    \end{pmatrix}\right) \\
                                                                 & = \left(I, \begin{pmatrix}
                                                                                      g_{11}a \\ 0
                                                                                  \end{pmatrix}\right)                                   \\
                                                                 & = \left(I, \begin{pmatrix}
                                                                                      ka' \\ 0
                                                                                  \end{pmatrix}\right)
    \end{align*}where the last equality comes from the RHS of Equation $\ref{eq:first-conjugacy}$. This now gives that
    $$g_{11} a = ka' \implies g_{11} a = \pm a'$$
    as $k = \pm 1$. But $a,a' > 0$ by the classification and so $g_{11} = k$ and $a = a'$.

    Next we inspect Equation $\ref{eq:second-conjugacy}$. The matrix part gives
    that $$ GBG^{-1} = B'^s .$$ Computing this gives
    \begin{align*}
        \begin{pmatrix}
            g_{11} & g_{12} \\
            0      & g_{22}
        \end{pmatrix}\begin{pmatrix}
                         1 & n \\
                         0 & 1
                     \end{pmatrix}\begin{pmatrix}
                                      g_{11} & -g_{12}g_{11}g_{22} \\
                                      0      & g_{22}
                                  \end{pmatrix} & = \begin{pmatrix}
                                                        1 & g_{11}g_{22}n \\
                                                        0 & 1
                                                    \end{pmatrix} \\
                                        & =\begin{pmatrix}
                                               1 & n's \\
                                               0 & 1
                                           \end{pmatrix}.
    \end{align*}
    Where the last equality comes again from the RHS of Equation $\ref{eq:second-conjugacy}$. Comparing matrix entries yields
    $g_{11}g_{22}n = n's$. Using that $g_{11} = \pm 1$,  $g_{22} = \pm 1$ and $s = \pm 1$ we have that $n = \pm n'$. But both $n,n' \in \mathbb{N}$ by the classification and so we must have $n = n'$.
    This now gives that $s = g_{11}g_{22}$. This was the linear part of Equation $\ref{eq:second-conjugacy}$ so we now analyze the translational component.

    Comparing the translational parts yields

    \begin{equation}
        \label{eq:vector-side}
        -GBG^{-1}\bar{h} + G\bar{b} + \bar{h} = s\bar{b}' + l\bar{a}'.
    \end{equation}
    We can factorise the LHS as
    $$
        -GBG^{-1}h + G\bar{b} + \bar{h}  = G(I- B)G^{-1}\bar{h} + G\bar{b}.
    $$
    Now we simplify
    $$
        G(I- B)G^{-1}\bar{h} = \begin{pmatrix}
            x \\ 0
        \end{pmatrix}
    $$
    where $x \in \mathbb{R}$ is arbritrary depending on $\bar{h}$. Now Equation \ref{eq:vector-side} becomes
    $$
        \begin{pmatrix}
            x \\ 0
        \end{pmatrix} + G \bar{b} = s\bar{b}' + l\bar{a}'.
    $$
    Expanding this we are left with
    $$
        \begin{pmatrix}
            x + g_{12}b \\
            g_{22}b
        \end{pmatrix} = \begin{pmatrix}
            la' \\
            sb'
        \end{pmatrix}.
    $$
    This gives $g_{22}b = sb'$ and again as $b, b'>0$ and $g_{22}, s = \pm 1$ we have $b = b'$ and now we are finished as we showed $a=a'$, $b=b'$ and $n=n'$.

\end{proof}

This completes the classification of the integral affine structures on the
torus $\mathbb{T}^2$. We note that lattices in the series
$\mathcal{A}_{a,b,c,d}$ induce a Euclidean structure on the torus - i.e. their
holonomy representation is a subgroup of $O(n) \rtimes \mathbb{R}^n$. This
implies that there is a Riemannian metric preserved by this structure. However,
the integral affine structures of $\mathcal{B}_{n,a,b}$ never induce a
Euclidean structure on the torus.

Next, we will look at the classification of integral affine structures over the
Klein bottle, hence completing the classification of all integral affine
structures over closed surfaces.

\newpage
\section{Classification of Integral Affine structures on the Klein Bottle}
\subsection{Completeness of integral affine structures on $K^2$}
We now wish to classify the integral affine structures on the Klein bottle. We
would hope to proceed as before by calculating all of the injective
homomorphisms form $\Gamma = \pi_1(K^2) \to \affz(\R^2)$ that induce a free and
proper action on $\mathbb{R}^2$. However, in order to do this we would need to
know that every integral affine structure on $K^2$ is complete. The following
lemma takes care of this fact.

\begin{lemma}
    Let $N$ be an integral affine manifold and suppose that there is a regular covering $M \to N$ by a compact orientable integral affine manifold whose fundamental group is nilpotent. Then any integral affine structure on $N$ is complete.
\end{lemma}
\begin{proof}
    We can pull back the integral affine structure on $N$ to $M$. A regular covering induces an injection of fundamental groups $\iota: \pi_1(M) \mapsto \pi_1(N)$. We can induce the holonomy representation on $M$ by $\hol \circ \iota$ where $\hol$ is the affine holonomy representation induced by the
    integral affine structure on $N$. As $M$ and $N$ share a common universal cover $X$, the integral affine structure on $N$ induces a developing map $\dev : X \mapsto \R^n$ and this induced holonomy representation will remain equivariant with respect to this developing map as the action of $\pi_1(\mathbb{T}^2)$ on the universal cover is a restriction of the action of $\pi_1(K^2)$.

    From Proposition $\ref{prop:dev-charts}$ this defines an integral affine
    structure on $M$ which has the same developing map as that of $N$. Therefore,
    if the integral affine structure on $N$ is incomplete then the integral affine
    structure on $M$ will also be. However, $M$ is a compact orientable integral
    affine manifold with nilpotent fundamental group and hence by Theorem
    \ref{thm:complete-parallel-volume} is complete. It follows that any integral
    affine structure on $N$ must also be complete.
\end{proof}

Throughout this section we let $\pi_1(K^2) = \Gamma$ have presentation $$\Gamma
    = \langle a,b \mid aba = b \rangle$$

and we note that $\langle a,b^2 \rangle \subset \Gamma$ is abelian. Indeed, the
oriented double cover of $K^2$ is the torus $\mathbb{T}^2$ and so there is an
injection of fundamental groups $\pi_1(\mathbb{T}^2) \hookrightarrow
    \pi_1(K^2)$. This abelian subgroup can be realised as the aforementioned
injection. As in the previous proposition, the induced integral affine
structure on $\mathbb{T}^2$ via pullback will have holonomy representation
$\hol \circ \iota$ where $\iota$ is the injection of fundamental groups and
$\hol$ is the holonomy representation of the affine structure on $\mathbb{T}^2$
- indeed this will be equivariant with respect to the same developing map on
their shared universal cover. This shows the following proposition.

\begin{prop} \label{prop:induced-structure}
    Let $K^2$ be equipped with an integral affine structure with holonomy representation
    $$\hol : \Gamma \to \affz(\R^2).$$
    Then if $\hol(a) = g \in \affz(\R^2)$ and $\hol(b) = h \in \affz(\R^2)$ then the induced holonomy representation on the torus $\mathbb{T}^2$
    will be given by
    $$\tilde{\hol} : \pi_1(\mathbb{T})^2 = \langle x,y \rangle \to \affz(\R^2)$$
    $$ x \mapsto g, \quad y \mapsto h^2$$
\end{prop}

\begin{remark}
    This implies that once we classify the injective homomorphisms from $\Gamma \to \affz(\R^2)$, we can easily check
    which integral affine structure they induce on the torus.
\end{remark}

\subsection{Determining the linear holonomy}
We now start are classification of the injective group morphisms from $\Gamma
    \to \affz(\R^2)$. We will denote also by $a = (A, \bar{x})$ and $b = (B,
    \bar{y})$ the images of our generators of $\Gamma$ under the holonomy
representation where $A,B \in GL(2, \Z)$ and $\bar{x}, \bar{y} \in
    \mathbb{R}^2$.

We will make use of the following in our classification

\begin{enumerate}
    \item We must have that $\det A = 1$ and $\det B^2 = 1$ as these are generators form
          $\pi_1(\mathbb{T}^2)$.
    \item Since $b$ reverses orientation we in fact have $\det B = -1$
    \item The action given by the composition $\pi_1(\mathbb{T}^2) \hookrightarrow \Gamma
              \to \affz(\R^2) $ must also act freely and properly as this defines the induced
          integral affine structure on $\mathbb{T}^2$ by Proposition
          \ref{prop:induced-structure}.
\end{enumerate}

We will now make use of Lemma \ref{lemma:traces-free-action} to get
restrictions on the matrix $B$. Namely, this lemma implies that $\text{Tr } B =
    0$. By Lemma 1 of \cite{torus-bundles}, our matrix $B$ must be conjugate to one
of $$ B_1 = \begin{pmatrix}
        1 & 0  \\
        0 & -1
    \end{pmatrix}\quad \text{or} \quad
    B_2=\begin{pmatrix}
        0 & 1 \\
        1 & 0
    \end{pmatrix}.
$$

Now we must find the matrices $A \in GL(2, \Z)$ such that $\det A = 1$,
$\text{Tr } A = 2$ which also satisfy the group relation $ABA = B$ where $B=
    B_1,B_2$. We perform this calculation now.

\begin{prop}
    Let $B = B_1$. Then we have the following choices for our matrix A.
    \begin{enumerate}
        \item $A = I$.
        \item $A_1 = \begin{pmatrix}
                      1 & n \\
                      0 & 1
                  \end{pmatrix}$.
        \item $A_2 = \begin{pmatrix}
                      1 & 0 \\
                      p & 1
                  \end{pmatrix}$.
    \end{enumerate}
\end{prop}

\begin{proof}
    Let $A = \begin{pmatrix}
            a & b \\
            c & d
        \end{pmatrix} \in GL(2, \Z)$. The trace condition implies that $d = 2-a$. We note that $B^{-1} = B$. Calculating the
    group condition $ABA = B$ we obtain
    \begin{align*}
        ABA & = \begin{pmatrix}
                    a^2 - bc & ab-bd    \\
                    ca -cd   & cb - d^2
                \end{pmatrix} \\
            & = \begin{pmatrix}
                    1 & 0  \\
                    0 & -1
                \end{pmatrix}.
    \end{align*}
    Comparing matrix entries we find
    $$a = d $$
    since $d = 2-a$ this implies $a=d=1$. Our matrix $A$ now has the form
    $$
        A = \begin{pmatrix}
            1 & b \\
            c & 1
        \end{pmatrix}.
    $$
    From $\det A = 1$ we find that either $b=0$ or $c = 0$ which completes the proof.
\end{proof}

Now we perform the same calculation for $B = B_2$.
\begin{prop}
    Let $B = B_2$. Then we have the following choices for our matrix A.
    \begin{enumerate}
        \item $A = I$.
        \item $A_3 = \begin{pmatrix}
                      1+n & n   \\
                      -n  & 1-n
                  \end{pmatrix}$.
        \item $A_4 = \begin{pmatrix}
                      1-n & n   \\
                      -n  & 1+n
                  \end{pmatrix}$.
    \end{enumerate}
\end{prop}
\begin{proof}
    Let $A= \begin{pmatrix}
            a & b \\
            c & d
        \end{pmatrix} \in GL(2, \Z) $. Again we must have that $d = 2-a$ by the trace condition. Calculating the group law $ABA =B$ yields
    \begin{align*}
        ABA & = \begin{pmatrix}
                    ab + ac  & b^2 + ad \\
                    ad + c^2 & bd + cd
                \end{pmatrix} \\
            & = \begin{pmatrix}
                    0 & 1 \\
                    1 & 0
                \end{pmatrix}.
    \end{align*}
    Comparing entries yields that $a(b+c) = 0$ and $d(b+c) = 0$. Hoiwever, the trace condition implies that $a$ and $d$ can't both be zero and so $b = -c$.
    Our matrix $A$ now has the form
    $$
        A = \begin{pmatrix}
            a  & b \\
            -b & d
        \end{pmatrix}
    $$
    Calculating the determinant condition and using $d = 2-a$ yields
    \begin{align*}
        b^2 & = 1-ad    \\
            & = (a-1)^2 \\
    \end{align*}
    and so
    $$ b = a-1\quad \text{or} \quad b = 1-a$$
    The first case yields
    $$
        A= \begin{pmatrix}
            n+1 & n   \\
            -n  & 1-n
        \end{pmatrix}
    $$
    where $n = a-1$ and the second case yields
    $$
        A= \begin{pmatrix}
            1-n & n   \\
            -n  & 1+n
        \end{pmatrix}
    $$
    where $n = 1-a$. The case where $n=0$ corresponds to the identity and this completes our proof.
\end{proof}

This gives the possible choices of $A$ and $B$ that satisfy the group law and
Lemma \ref{lemma:traces-free-action}. We must now deduce the translational
component in each case and check that the induced action is free and proper.

\subsection{Translational component and freeness of the action}
To find the translational components we impose the group law $aba = b$. When
checking freeness, we can often use Proposition \ref{prop:induced-structure} to
take shortcuts as we know that $a,b^2$ must generate one of the integral affine
structures on the torus and hence act freely.

We will now compute the translational components and check whether the action
is free.

\begin{lemma}
    Let $B = B_1$. Then the possible choices of $A$ that induce an integral affine structure on $K^2$ are
    $A=I, A_1$.
\end{lemma}
\begin{proof}
    \begin{case}[$A=I$, $B=B_1$]
        Let $a=(I, \bar{x})$ and $b = (B_1, \bar{b})$ with $\bar{x} = \begin{pmatrix}
                x_1 \\ x_2
            \end{pmatrix}$ and $\bar{b} = \begin{pmatrix}
                y_1 \\ y_2
            \end{pmatrix}$ be our generators in $\affz(\R^2)$. They must satisfy the relation $aba = b$. Calculating this we find
        \begin{align*}
            (B, B\bar{x} +\bar{x} + \bar{b}) = (B, \bar{b})
        \end{align*}
        The translational component gives $B\bar{x} = - \bar{x}$ and so $\bar{x}$ lives in the $-1$-eigenspace. This implies that $x_1=0$. We are free to conjugate both our generators as this preserves the integral affine structure
        and so we will do so by $g= \left(I, \begin{pmatrix}
                    0 \\ \frac{y_2}{2}
                \end{pmatrix}\right)$ Conjugating $a$ leaves it unchanged so we calculate

        $$
            g^{-1}\circ b \circ g = \left(B, \begin{pmatrix}
                    y_1 \\ 0
                \end{pmatrix}\right)
        $$
        and so we may assume that $y_2 =0$ by performing this conjugation. Our generators now look like
        $$a = \left(I, \begin{pmatrix}
                    0 \\ x_2
                \end{pmatrix}\right), \quad b = \left(B, \begin{pmatrix}
                    y_1 \\ 0
                \end{pmatrix}\right)$$
        Since $B$ alone clearly does not act freely we can't have that $y_1 =0$ and if $x_2=0$ then our lattice is one-dimensional and hence the holonomy representation is not injective.
        Therefore, $x_2, y_1 \neq 0$. Furthermore, by changing generators $a \mapsto a^{-1}$ and $b \mapsto b^{-1}$ we can assume $x_2, y_1 >0$. It remains to check that this action is free.
        We must check that $a^qb^n$ acts freely for all $q,n \in \Z$ not both zero. However, we now note that as $B^{2k} = I$ for $k\in \Z$, by the torus classification of Theorem \ref{thm:torus-series-a},
        we have already shown that $a^qb^{2k}$ acts freely. We therefore must check the case of $a^qb^{2k+1}$. We compute that
        $$
            a^q = \left(I, \begin{pmatrix}
                    0 \\ qx_2
                \end{pmatrix} \right)
            , \quad b^{2k+1} = \left(B, \begin{pmatrix}
                    (2k+1)y_1 \\ 0
                \end{pmatrix} \right)$$
        For the action to be free we need $a^qb^{2k+1}\circ z = z$ to have no solutions for all $z \in \R^2$. We compute this action
        \begin{align*}
            a^qb^{2k+1}\circ z & = B \begin{pmatrix}
                                         z_1 \\ z_2
                                     \end{pmatrix} + \begin{pmatrix}
                                                         (2k+1)y_1 \\ x_2
                                                     \end{pmatrix} \\
                               & = \begin{pmatrix}
                                       (2k+1)y_1 + z_1 \\ x_2 - z_2
                                   \end{pmatrix}
        \end{align*}
        and so, for there to be a fixed point we need
        $$(2k+1)y_1 + z_1 = z_1$$
        which gives $(2k+1)y_1 = 0$. But now $y_1>0$ and $2k+1 = 0 \implies k = -\frac{1}{2} \notin \Z$ and so this action can have no fixed points. The induced integral affine structure on the torus is given by
        $$
            a = \left(1, \begin{pmatrix}
                    0 \\ x_2
                \end{pmatrix} \right), \quad b^2 = \left(1, \begin{pmatrix}
                    2y_1 \\ 0
                \end{pmatrix} \right)
        $$
    \end{case}
    \begin{case}[$A=A_1$, $B=B_1$]
        Let $a = \left( A_1, \begin{pmatrix}
                    x_1 \\ x_2
                \end{pmatrix} \right)$ and $b = \left( B, \begin{pmatrix}
                    y_1 \\ y_2
                \end{pmatrix} \right)$
        As $A_1^2 \neq I$ and $B^2 = I$, by the torus classification we again obtain that $x_1 = 0$. We calculate the group law $aba = b$.
        This gives us
        \begin{align*}
            aba & = \left(B, \begin{pmatrix}
                                 y_1 + n(y_2-x_2) \\
                                 y_2
                             \end{pmatrix}\right)   \\
                & = \left( B, \begin{pmatrix}
                                  y_1 \\ y_2
                              \end{pmatrix} \right)
        \end{align*}
        which imposes the restriction $n(y_2 - x_2) = 0$ and since $n \neq 0$ this implies $y_2 = x_2$. We can also assume that $x_2>0$ by conjugating the generators by $(_I, \bar{0})$ and that $y_1>0$ by changing generators $b \mapsto b^{-1}$ if necessary.
        Now we check freeness of the action.

        We calculate that $$ a^q = \left( \begin{pmatrix}
                    1 & qn \\ 0 & 1
                \end{pmatrix}, \begin{pmatrix}
                    \frac{nx_2 q(q-1)}{2} \\ qx_2
                \end{pmatrix} \right)
        $$
        and
        $$b^k = \begin{cases*}
                \left( B, \begin{pmatrix}
                                  ky_1 \\ x_2
                              \end{pmatrix} \right) \text{, k odd } \\
                \left( I, \begin{pmatrix}
                                  ky_1 \\ 0
                              \end{pmatrix} \right)\text{, k even }
            \end{cases*}$$
        Since the case of $k$ being even is covered by the torus classification we must only check freeness for $k$ being odd.

        For $k$ odd we compute
        \begin{align*}
            a^qb^k = \left( \begin{pmatrix}
                                1 & -qn \\
                                0 & -1
                            \end{pmatrix}, \begin{pmatrix}
                                               \dfrac{nx_2q(q-1)}{2} + qnx_2 + ky_1 \\
                                               (q+1)x_2
                                           \end{pmatrix} \right)
        \end{align*}
        calculating this action on $z = \begin{pmatrix}
                z_1 \\ z_2
            \end{pmatrix}$ yields

        $$
            (a^qb^k) \circ z = \begin{pmatrix}
                \dfrac{nx_2q(q-1)}{2} + qnx_2 + ky_1 + z_1 - qnz_2 \\
                (q+1)x_2 - z_2
            \end{pmatrix}
        $$
        Solving the fixed point equation, the bottom entry of the vector yields
        \begin{equation}
            \label{eq:fixed-point-case-two}
            z_2 = \dfrac{q+1}{2}x_2
        \end{equation}
        Solving the fixed point equation in the top component yields
        $$
            z_2 = \dfrac{x_2(q-1)}{2}x_2 + \dfrac{ky_1}{qn}
        $$
        Substituting Equation \ref{eq:fixed-point-case-two} for $z_2$ yields $ky_1 = 0$ which implies $k=0$ as $y_1>0$.
        This now implies that $a^q$ would need to have a fixed point but one checks

        $$a^q \circ z = \begin{pmatrix}
                * \\
                z_2 + qx_2
            \end{pmatrix}$$
        and so we require that $qx_2 =0$ but $x_2>0$ and hence the only element with a fixed point occurs when $k=q=0$ which is the identity and hence the action is free.

        The induced integral affine structure on $\mathbb{T}^2$ is given by

        $$
            a = \left(A_1, \begin{pmatrix}
                    0 \\ x_2
                \end{pmatrix} \right), \quad b^2 = \left(I, \begin{pmatrix}
                    2y_1 \\ 0
                \end{pmatrix} \right)
        $$
        with $x_2, y_1 >0$
    \end{case}
    \begin{case}[$A=A_2$, $B=B_1$]
        Let $a = \left( A_1, \begin{pmatrix}
                    x_1 \\ x_2
                \end{pmatrix} \right)$ and $b = \left( B, \begin{pmatrix}
                    y_1 \\ y_2
                \end{pmatrix} \right)$.
        In this case the action will fail to be free. Again, by considering the induced structure on the torus it follows that $x_1=0$. We calculate $aba$ again which yields
        \begin{align*}
            aba & = \left(B, \begin{pmatrix}
                                 y_1 \\ py_1 + y_2
                             \end{pmatrix}\right) \\
                & = \left(B, \begin{pmatrix}
                                 y_1 \\ y_2
                             \end{pmatrix}\right)
        \end{align*}

        and so we must have $py_1 = 0$ and since $p>0$ we have $y_1=0$. But now we
        calculate $b^2$ as $$ b^2 = \left(I, \begin{pmatrix}
                    2y_1 \\ 0
                \end{pmatrix}\right)
        $$
        but since $y_1=0$ this cannot induce an integral affine structure on the torus hence the action is not free.
    \end{case}
\end{proof}
This completes the case where $B = B_1$. Next, we tackle the case of $B = B_2$.
\begin{lemma}
    Let $B = B_2$. Then the possible choices of $A$ that induce an integral affine structure on $K^2$ are
    $A=I, A_4$.
\end{lemma}
\begin{proof}
    \begin{case}[$A=I$, $B=B_2$]
        We calculate the group law $aba = b$ which yields
        \begin{align*}
            aba & = \left(B, \begin{pmatrix}
                                 x_1 +x_2 +y_1 \\ x_2 + y_2 + x_1
                             \end{pmatrix}\right) \\
                & = \left(B, \begin{pmatrix}
                                 y_1 \\ y_2
                             \end{pmatrix}\right)
        \end{align*}
        This implies that $x_1 = -x_2$. We conjugate $a,b$ by
        $$
            c = \left(-I, \begin{pmatrix}
                    -\frac{y_2}{2} \\
                    -\frac{y_2}{2}
                \end{pmatrix} \right)
        $$
        which gives use the new generators
        $$
            a = \left(I , \begin{pmatrix}
                    -x_1 \\ x_1
                \end{pmatrix}\right), \quad b =\left(B , \begin{pmatrix}
                    -y_2 -y_1 \\ 0
                \end{pmatrix}\right)
        $$
        and by switching generators $a \mapsto a^{-1}$ if necessary we may assume that $x_1<0$ so that we will switch the signs on the translational component of $a$.
        We will do the same for the translational component of $b$ and assume $y_2 = 0$ and $y_1>0$ after relabelling $-y_2 -y_1 \to y_1$. This finishes our generating set which we write as
        $$
            a = \left(I , \begin{pmatrix}
                    x_1 \\ -x_1
                \end{pmatrix}\right), \quad b =\left(B , \begin{pmatrix}
                    y_1 \\ 0
                \end{pmatrix}\right)
        $$
        with $x_1, y_1>0$. We now check freeness of the action. Again since $B^{2k}= I$, the even case is covered by the torus classification and so
        we will check freeness of words $a^qb^k$ with $k$ being odd. We calculate that
        \begin{align*}
             & a^q = \left(I, \begin{pmatrix}
                                      qx_1 \\ - qx_1
                                  \end{pmatrix}\right)                                  \\
             & b^k = \left(B, \begin{pmatrix}
                                      \frac{k+1}{2}y_1 \\ \frac{k-1}{2}y_1
                                  \end{pmatrix}\right) \text{, for k odd}               \\
             & a^qb^k = \left(B, \begin{pmatrix}
                                         \frac{k+1}{2}y_1  + qx_1 \\ \frac{k-1}{2}y_1 - qx_1
                                     \end{pmatrix}\right) \text{, for k odd.}
        \end{align*}
        Now let $z = \begin{pmatrix}
                z_1 \\ z_2
            \end{pmatrix} \in \R^2$ and consider the fixed point equation $(a^qb^k)\circ z = z$. This gives

        $$
            \begin{pmatrix}
                \frac{k+1}{2}y_1  + qx_1 + z_2 \\
                \frac{k-1}{2}y_1  - qx_1 + z_1
            \end{pmatrix} = \begin{pmatrix}
                z_1 \\ z_2
            \end{pmatrix}
        $$
        Solving this gives $\frac{k+1}{2}y_1 + \frac{k-1}{2}y_1 = 0 \implies ky_1 = 0$. But since $y_1>0$ we must have $k=0$. So if the action is not free we must have that $a^q$ acts non-freely for some $q\neq 0$
        but this action is by translations and hence is free and so $A=I, B=B_2$ does induce an integral affine structure. The induced structure on the torus is given by
        $$
            a = \left(I , \begin{pmatrix}
                    x_1 \\ -x_1
                \end{pmatrix}\right), \quad b^2 =\left(I , \begin{pmatrix}
                    y_1 \\ y_1
                \end{pmatrix}\right)
        $$
        $x_1, y_1 > 0$.
    \end{case}
    \begin{case}[$A=A_3$, $B=B_2$]
        We show that this cannot induce an integral affine structure on the torus. We conjugate $A$ and $B$ by
        $$
            G = \begin{pmatrix}
                1 & 0 \\
                1 & 1
            \end{pmatrix}
        $$
        which puts them in the form
        $$
            A = \begin{pmatrix}
                1 & n \\
                0 & 1
            \end{pmatrix}, \quad B = \begin{pmatrix}
                -1 & 1 \\
                0  & 1
            \end{pmatrix}
        $$
        and so we have $B^2 = I$. Hence, the integral affine structure on $\mathbb{T}^2$ should belong to the second family. However calculating
        $$
            b^2 = \left(I, \begin{pmatrix}
                    y_2 \\ 2y_2
                \end{pmatrix}\right)
        $$
        we require that $2y_2 = 0 \implies y_2 = 0$ but then $b^2$ does not act freely and hence does not induce an integral affine structure on the torus. Therefore, this case cannot induce
        an integral affine structure on $K^2$.
    \end{case}
    \begin{case}[$A=A_4$, $B=B_2$]
        Again, we conjugate $A,B$ by the matrix
        $$
            G = \begin{pmatrix}
                1  & 0 \\
                -1 & 1
            \end{pmatrix}
        $$
        and assume that
        $$
            a = \left(A = \begin{pmatrix}
                    1 & n \\
                    0 & 1
                \end{pmatrix} , \begin{pmatrix}
                    x_1 \\ x_2
                \end{pmatrix}\right), \quad b =\left(B = \begin{pmatrix}
                    1 & 1  \\
                    0 & -1
                \end{pmatrix} , \begin{pmatrix}
                    y_1 \\ y_2
                \end{pmatrix}\right)
        $$
        Using that $B^2 = I$ and $a,b^2$ must induce an integral affine structure on the torus it follows that $x_1 = 0$ since the induced structure belongs to the second family.
        Now, calculating the group law $aba = b$ yields
        $$
            \left(\begin{pmatrix}
                    1 & 1  \\
                    0 & -1
                \end{pmatrix} , \begin{pmatrix}
                    y_1 + ny_2 + x_2(1-n) \\ y_2
                \end{pmatrix}\right) = \left(B , \begin{pmatrix}
                    y_1 \\ y_2
                \end{pmatrix}\right)
        $$
        and so $ny_2 + (1-n)x_2 = 0 \implies y_2 = \frac{n-1}{n}x_2$. Our generators are therefore
        $$
            a = \left(\begin{pmatrix}
                    1 & n \\
                    0 & 1
                \end{pmatrix} , \begin{pmatrix}
                    0 \\ x_2
                \end{pmatrix}\right), \quad b =\left(\begin{pmatrix}
                    1 & 1  \\
                    0 & -1
                \end{pmatrix} , \begin{pmatrix}
                    y_1 \\ \frac{n-1}{n}x_2
                \end{pmatrix}\right)
        $$
        and the induced structure on the torus is
        $$
            a = \left(\begin{pmatrix}
                    1 & n \\
                    0 & 1
                \end{pmatrix} , \begin{pmatrix}
                    0 \\ x_2
                \end{pmatrix}\right), \quad b^2 =\left(\begin{pmatrix}
                    1 & 0 \\
                    0 & 1
                \end{pmatrix} , \begin{pmatrix}
                    2y_1 + \frac{n-1}{n}x_2 \\ 0
                \end{pmatrix}\right)
        $$
        with $y_1> 0$ and so a necessary condition is $2y_1 \neq -\frac{n-1}{n}x_2$. In fact, we shall now show that this is sufficient. We calculate that
        $$
            a^k = \left(\begin{pmatrix}
                    1 & kn \\
                    0 & 1
                \end{pmatrix} , \begin{pmatrix}
                    \frac{nk(k-1)}{2}x_2 \\ kx_2
                \end{pmatrix}\right), \quad b^q =\left(B , \begin{pmatrix}
                    qy_1 + \frac{q-1}{2}y_2 \\ y_2
                \end{pmatrix}\right) \text{, for q odd}
        $$
        and
        $$
            a^kb^q = \left(\begin{pmatrix}
                    1 & 1-kn \\
                    0 & -1
                \end{pmatrix} , \begin{pmatrix}
                    qy_1 + \frac{q-1}{2}y_2 + kny_2 + \frac{nk(k-1)}{2}x_2 \\ y_2 + kx_2
                \end{pmatrix}\right)
        $$
        Letting $v_1 := qy_1 + \frac{q-1}{2}y_2 + kny_2 + \frac{nk(k-1)}{2}x_2$ and $v_2 := y_2 + kx_2$ and $z = \begin{pmatrix}
                z_1 \\ z_2
            \end{pmatrix} \in \R^2$, the fixed point equation yields
        $$
            \begin{pmatrix}
                (1-kn)z_2 \\
                -2z_2
            \end{pmatrix} = - \begin{pmatrix}
                v_1 \\
                v_2
            \end{pmatrix}
        $$
        This yields the equation $(1-kn)v_2 + 2v_1 = 0$. Solving this equation and using that $y_2 = \frac{n-1}{n}x_2$  we obtain
        $$
            q(2y_1 +y_2) =0
        $$
        It is quick to check that $a^k$ acts freely so if the action is not free then $q \neq 0$, and this gives the requirement
        $$
            2y_1 \neq -\frac{n-1}{n}x_2
        $$
        and so this condition is sufficient and necessary.
    \end{case}
\end{proof}
This completes the classification of the integral affine structures on the torus. We summarize the results of the previous two lemmas into the following theorem.
\begin{theorem}
    The possible integral integral affine structures on the Klein bottle are
    \begin{enumerate}
        \item $a = \left(I, \begin{pmatrix}
                      0 \\ x_2
                  \end{pmatrix}\right), \quad b = \left(\begin{pmatrix}
                      1 & 0  \\
                      0 & -1
                  \end{pmatrix}, \begin{pmatrix}
                      y_1 \\ 0
                  \end{pmatrix}\right)$, with $x_2, y_1>0$.
        \item $a = \left(\begin{pmatrix}
                      1 & n \\
                      0 & 1
                  \end{pmatrix}, \begin{pmatrix}
                      0 \\ x_2
                  \end{pmatrix}\right), \quad b = \left(\begin{pmatrix}
                      1 & 0  \\
                      0 & -1
                  \end{pmatrix}, \begin{pmatrix}
                      y_1 \\ x_2
                  \end{pmatrix}\right)$, with $x_2, y_1>0$.
        \item $a = \left(I, \begin{pmatrix}
                      x_1 \\ -x_1
                  \end{pmatrix}\right), \quad b = \left(\begin{pmatrix}
                      0 & 1 \\
                      1 & 0
                  \end{pmatrix}, \begin{pmatrix}
                      y_1 \\ 0
                  \end{pmatrix}\right)$, with $x_1, y_1>0$.
        \item $a = \left(\begin{pmatrix}
                          1 & n \\
                          0 & 1
                      \end{pmatrix} , \begin{pmatrix}
                          0 \\ x_2
                      \end{pmatrix}\right), \quad b =\left(\begin{pmatrix}
                          1 & 1  \\
                          0 & -1
                      \end{pmatrix} , \begin{pmatrix}
                          y_1 \\ \frac{n-1}{n}x_2
                      \end{pmatrix}\right)$, with $y_1> 0$ and \\
              $2y_1 \neq -\frac{n-1}{n}x_2$.
    \end{enumerate}
\end{theorem}

We note that families 1) and 3) induce a Euclidean structure on the Klein
bottle, e.g. they preserves some Riemannian metric as we have that the group
they generate is a subset of $O(n) \rtimes \R^2$.

We now wish to classify the corresponding Lagrangian fibrations for each of the
integral affine structures on the two torus and the Klein bottle. Below are
some pictures of the different families of integral affine structures on both
the torus and the Klein bottle. (TODO: Lattices for the Klein Bottle)

\subsection{Fundamental domains for $\mathbb{T}^2$ lattices.}
\begin{figure}[H]
    \centering
    \incfig[0.7]{secondlatticeguy}
    \caption{Fundamental domain for lattice $\mathcal{A}_{a,b,c,d}$}
\end{figure}

\begin{figure}[H]
    \centering
    \incfig[0.7]{firstlattice}
    \caption{Fundamental domain for lattice $\mathcal{B}_{1,1,1}$}
\end{figure}